% Options for packages loaded elsewhere
% Options for packages loaded elsewhere
\PassOptionsToPackage{unicode}{hyperref}
\PassOptionsToPackage{hyphens}{url}
%
\documentclass[
  11pt,
  ignorenonframetext,
  aspectratio=169,
]{beamer}
\newif\ifbibliography
\usepackage{pgfpages}
\setbeamertemplate{caption}[numbered]
\setbeamertemplate{caption label separator}{: }
\setbeamercolor{caption name}{fg=normal text.fg}
\beamertemplatenavigationsymbolsempty
% remove section numbering
\setbeamertemplate{part page}{
  \centering
  \begin{beamercolorbox}[sep=16pt,center]{part title}
    \usebeamerfont{part title}\insertpart\par
  \end{beamercolorbox}
}
\setbeamertemplate{section page}{
  \centering
  \begin{beamercolorbox}[sep=12pt,center]{section title}
    \usebeamerfont{section title}\insertsection\par
  \end{beamercolorbox}
}
\setbeamertemplate{subsection page}{
  \centering
  \begin{beamercolorbox}[sep=8pt,center]{subsection title}
    \usebeamerfont{subsection title}\insertsubsection\par
  \end{beamercolorbox}
}
% Prevent slide breaks in the middle of a paragraph
\widowpenalties 1 10000
\raggedbottom
\AtBeginPart{
  \frame{\partpage}
}
\AtBeginSection{
  \ifbibliography
  \else
    \frame{\sectionpage}
  \fi
}
\AtBeginSubsection{
  \frame{\subsectionpage}
}
\usepackage{iftex}
\ifPDFTeX
  \usepackage[T1]{fontenc}
  \usepackage[utf8]{inputenc}
  \usepackage{textcomp} % provide euro and other symbols
\else % if luatex or xetex
  \usepackage{unicode-math} % this also loads fontspec
  \defaultfontfeatures{Scale=MatchLowercase}
  \defaultfontfeatures[\rmfamily]{Ligatures=TeX,Scale=1}
\fi
\usepackage{lmodern}

\ifPDFTeX\else
  % xetex/luatex font selection
\fi
% Use upquote if available, for straight quotes in verbatim environments
\IfFileExists{upquote.sty}{\usepackage{upquote}}{}
\IfFileExists{microtype.sty}{% use microtype if available
  \usepackage[]{microtype}
  \UseMicrotypeSet[protrusion]{basicmath} % disable protrusion for tt fonts
}{}
\makeatletter
\@ifundefined{KOMAClassName}{% if non-KOMA class
  \IfFileExists{parskip.sty}{%
    \usepackage{parskip}
  }{% else
    \setlength{\parindent}{0pt}
    \setlength{\parskip}{6pt plus 2pt minus 1pt}}
}{% if KOMA class
  \KOMAoptions{parskip=half}}
\makeatother

\usepackage{color}
\usepackage{fancyvrb}
\newcommand{\VerbBar}{|}
\newcommand{\VERB}{\Verb[commandchars=\\\{\}]}
\DefineVerbatimEnvironment{Highlighting}{Verbatim}{commandchars=\\\{\}}
% Add ',fontsize=\small' for more characters per line
\usepackage{framed}
\definecolor{shadecolor}{RGB}{241,243,245}
\newenvironment{Shaded}{\begin{snugshade}}{\end{snugshade}}
\newcommand{\AlertTok}[1]{\textcolor[rgb]{0.68,0.00,0.00}{#1}}
\newcommand{\AnnotationTok}[1]{\textcolor[rgb]{0.37,0.37,0.37}{#1}}
\newcommand{\AttributeTok}[1]{\textcolor[rgb]{0.40,0.45,0.13}{#1}}
\newcommand{\BaseNTok}[1]{\textcolor[rgb]{0.68,0.00,0.00}{#1}}
\newcommand{\BuiltInTok}[1]{\textcolor[rgb]{0.00,0.23,0.31}{#1}}
\newcommand{\CharTok}[1]{\textcolor[rgb]{0.13,0.47,0.30}{#1}}
\newcommand{\CommentTok}[1]{\textcolor[rgb]{0.37,0.37,0.37}{#1}}
\newcommand{\CommentVarTok}[1]{\textcolor[rgb]{0.37,0.37,0.37}{\textit{#1}}}
\newcommand{\ConstantTok}[1]{\textcolor[rgb]{0.56,0.35,0.01}{#1}}
\newcommand{\ControlFlowTok}[1]{\textcolor[rgb]{0.00,0.23,0.31}{\textbf{#1}}}
\newcommand{\DataTypeTok}[1]{\textcolor[rgb]{0.68,0.00,0.00}{#1}}
\newcommand{\DecValTok}[1]{\textcolor[rgb]{0.68,0.00,0.00}{#1}}
\newcommand{\DocumentationTok}[1]{\textcolor[rgb]{0.37,0.37,0.37}{\textit{#1}}}
\newcommand{\ErrorTok}[1]{\textcolor[rgb]{0.68,0.00,0.00}{#1}}
\newcommand{\ExtensionTok}[1]{\textcolor[rgb]{0.00,0.23,0.31}{#1}}
\newcommand{\FloatTok}[1]{\textcolor[rgb]{0.68,0.00,0.00}{#1}}
\newcommand{\FunctionTok}[1]{\textcolor[rgb]{0.28,0.35,0.67}{#1}}
\newcommand{\ImportTok}[1]{\textcolor[rgb]{0.00,0.46,0.62}{#1}}
\newcommand{\InformationTok}[1]{\textcolor[rgb]{0.37,0.37,0.37}{#1}}
\newcommand{\KeywordTok}[1]{\textcolor[rgb]{0.00,0.23,0.31}{\textbf{#1}}}
\newcommand{\NormalTok}[1]{\textcolor[rgb]{0.00,0.23,0.31}{#1}}
\newcommand{\OperatorTok}[1]{\textcolor[rgb]{0.37,0.37,0.37}{#1}}
\newcommand{\OtherTok}[1]{\textcolor[rgb]{0.00,0.23,0.31}{#1}}
\newcommand{\PreprocessorTok}[1]{\textcolor[rgb]{0.68,0.00,0.00}{#1}}
\newcommand{\RegionMarkerTok}[1]{\textcolor[rgb]{0.00,0.23,0.31}{#1}}
\newcommand{\SpecialCharTok}[1]{\textcolor[rgb]{0.37,0.37,0.37}{#1}}
\newcommand{\SpecialStringTok}[1]{\textcolor[rgb]{0.13,0.47,0.30}{#1}}
\newcommand{\StringTok}[1]{\textcolor[rgb]{0.13,0.47,0.30}{#1}}
\newcommand{\VariableTok}[1]{\textcolor[rgb]{0.07,0.07,0.07}{#1}}
\newcommand{\VerbatimStringTok}[1]{\textcolor[rgb]{0.13,0.47,0.30}{#1}}
\newcommand{\WarningTok}[1]{\textcolor[rgb]{0.37,0.37,0.37}{\textit{#1}}}

\usepackage{longtable,booktabs,array}
\newcounter{none} % for unnumbered tables
\usepackage{calc} % for calculating minipage widths
\usepackage{caption}
% Make caption package work with longtable
\makeatletter
\def\fnum@table{\tablename~\thetable}
\makeatother
\usepackage{graphicx}
\makeatletter
\newsavebox\pandoc@box
\newcommand*\pandocbounded[1]{% scales image to fit in text height/width
  \sbox\pandoc@box{#1}%
  \Gscale@div\@tempa{\textheight}{\dimexpr\ht\pandoc@box+\dp\pandoc@box\relax}%
  \Gscale@div\@tempb{\linewidth}{\wd\pandoc@box}%
  \ifdim\@tempb\p@<\@tempa\p@\let\@tempa\@tempb\fi% select the smaller of both
  \ifdim\@tempa\p@<\p@\scalebox{\@tempa}{\usebox\pandoc@box}%
  \else\usebox{\pandoc@box}%
  \fi%
}
% Set default figure placement to htbp
\def\fps@figure{htbp}
\makeatother





\setlength{\emergencystretch}{3em} % prevent overfull lines

\providecommand{\tightlist}{%
  \setlength{\itemsep}{0pt}\setlength{\parskip}{0pt}}



 


\setbeamertemplate{navigation symbols}{}
\setbeamertemplate{footline}[frame number]
\setbeamertemplate{itemize items}[circle]
\setbeamertemplate{itemize subitem}[triangle]
\setbeamercolor{frametitle}{fg=black}
\setbeamerfont{frametitle}{series=\bfseries}
\usepackage{amsmath}
\newcommand{\indep}{\perp\!\!\!\perp}
\usepackage{booktabs}
\usepackage{array}
\newcolumntype{R}[1]{>{\raggedright\arraybackslash}p{#1}}
\usepackage{fancyvrb}
\fvset{fontsize=\scriptsize}
\makeatletter
\def\verbatim@font{\scriptsize\ttfamily}
\makeatother
\makeatletter
\@ifpackageloaded{caption}{}{\usepackage{caption}}
\AtBeginDocument{%
\ifdefined\contentsname
  \renewcommand*\contentsname{Table of contents}
\else
  \newcommand\contentsname{Table of contents}
\fi
\ifdefined\listfigurename
  \renewcommand*\listfigurename{List of Figures}
\else
  \newcommand\listfigurename{List of Figures}
\fi
\ifdefined\listtablename
  \renewcommand*\listtablename{List of Tables}
\else
  \newcommand\listtablename{List of Tables}
\fi
\ifdefined\figurename
  \renewcommand*\figurename{Figure}
\else
  \newcommand\figurename{Figure}
\fi
\ifdefined\tablename
  \renewcommand*\tablename{Table}
\else
  \newcommand\tablename{Table}
\fi
}
\@ifpackageloaded{float}{}{\usepackage{float}}
\floatstyle{ruled}
\@ifundefined{c@chapter}{\newfloat{codelisting}{h}{lop}}{\newfloat{codelisting}{h}{lop}[chapter]}
\floatname{codelisting}{Listing}
\newcommand*\listoflistings{\listof{codelisting}{List of Listings}}
\makeatother
\makeatletter
\makeatother
\makeatletter
\@ifpackageloaded{caption}{}{\usepackage{caption}}
\@ifpackageloaded{subcaption}{}{\usepackage{subcaption}}
\makeatother

\usepackage{bookmark}
\IfFileExists{xurl.sty}{\usepackage{xurl}}{} % add URL line breaks if available
\urlstyle{same}
\hypersetup{
  pdftitle={Difference-in-Differences: Complex Designs},
  pdfauthor={Professor Ji-Woong Chung},
  hidelinks,
  pdfcreator={LaTeX via pandoc}}


\title{\texorpdfstring{Difference-in-Differences: Complex
Designs}{Difference-in-Differences: Complex Designs}}
\subtitle{\texorpdfstring{BUSS975: Causal Inference in Financial
Research}{BUSS975: Causal Inference in Financial Research}}
\author{Professor Ji-Woong Chung}
\date{}
\institute{Korea University}

\begin{document}
\frame{\titlepage}


\begin{frame}{Where we left off}
\protect\phantomsection\label{where-we-left-off}
\begin{itemize}
\tightlist
\item
  Lecture 7a: difference-in-differences is \textbf{four averages and
  three subtractions}. It identifies the ATT if the untreated outcomes
  of the two groups would have moved in parallel, and if the baseline is
  untreated.
\item
  Three regressions computed the same number: the interaction, two-way
  fixed effects, the long difference. ``DiD'' and ``TWFE'' were
  synonyms.
\end{itemize}

\vspace{0.4em}

\begin{itemize}
\tightlist
\item
  This lecture is about what happens when the design stops being a
  single \(2 \times 2\):

  \begin{itemize}
  \tightlist
  \item
    the sample changes under your feet;
  \item
    the trends are not parallel, but a second comparison group might be;
  \item
    the groups differ in ways you can measure;
  \item
    and above all, \textbf{treatment arrives at different dates}, and
    two-way fixed effects quietly turns into a weighted average of
    comparisons you would never have chosen.
  \end{itemize}
\item
  Source: Cunningham (2025), \emph{The Remix}, chapter 10. The
  simulation data are his; the castle-doctrine data are from 7a.
\end{itemize}
\end{frame}

\begin{frame}{The plan}
\protect\phantomsection\label{the-plan}
\begin{enumerate}
\tightlist
\item
  \textbf{Who is in the sample.} Attrition in panels; composition in
  repeated cross-sections.
\item
  \textbf{Triple differences.} A design, not a falsification.
\item
  \textbf{Covariates.} Conditional parallel trends and four ways to use
  it: weighting, imputation, both, or the regression everyone runs.
\item
  \textbf{Differential timing and TWFE.} The Bacon decomposition: what
  the coefficient averages, with what weights, and how it can get the
  sign wrong.
\item
  \textbf{Estimators that refuse the forbidden comparison.}
  Callaway--Sant'Anna, Sun--Abraham, de Chaisemartin--D'Haultfoeuille,
  imputation.
\item
  \textbf{Reading an event study.} Long differences, short gaps, honest
  bounds.
\item
  \textbf{A checklist} for designing your own.
\end{enumerate}

\vspace{0.4em}

\begin{itemize}
\tightlist
\item
  One running question: \textbf{which \(2 \times 2\)'s are inside my
  estimate, and what weight does each carry?}
\end{itemize}
\end{frame}

\section{1. Who Is in the Sample}\label{who-is-in-the-sample}

\begin{frame}{Imbalanced panels}
\protect\phantomsection\label{imbalanced-panels}
\begin{itemize}
\tightlist
\item
  A panel follows the same units. When some drop out or reappear, it is
  \textbf{imbalanced}. Does that break parallel trends?
\item
  \textbf{Not necessarily.} Parallel trends is a statement about the
  \emph{average change} in \(Y^0\). If units go missing \textbf{at
  random} --- missingness unrelated to \(Y^0\) --- the average is
  unaffected (Wooldridge 2010).
\item
  Weaker than full missing-at-random: it is enough that missingness is
  independent of the \emph{untreated} outcome given covariates, \[
  Y^0 \;\indep\; M \mid X \qquad\Longrightarrow\qquad \mathbb{E}[Y^0 \mid M=1, X] = \mathbb{E}[Y^0 \mid M=0, X].
  \]
\item
  Maria leaves because training got her promoted: missing because of her
  \(Y^1\), not her \(Y^0\). That does not bias the control trend, and it
  allows imputing a control unit's missing \emph{observed} outcome from
  other untreated units. A missing-data problem, not an identification
  problem.
\item
  Know \textbf{why} units are missing before you discard the design or
  the data. Chained DiD (Bellego, Benatia and Dortet-Bernadet 2024)
  links adjacent-period changes instead of throwing data away, at the
  cost of its own assumptions.
\end{itemize}
\end{frame}

\begin{frame}{Repeated cross-sections: Napster and composition}
\protect\phantomsection\label{repeated-cross-sections-napster-and-composition}
Hong (2013): did Napster (1999--2001) reduce music spending? No
geographic variation in access, so the comparison is \textbf{internet
users} vs.~\textbf{non-users} in the Consumer Expenditure Survey --- a
fresh random sample every year.

\begin{center}\small
\begin{tabular}{lcccccccc}
\toprule
 & \multicolumn{2}{c}{1997} & \multicolumn{2}{c}{1998} & \multicolumn{2}{c}{1999} & \multicolumn{2}{c}{2000} \\
 & User & Non & User & Non & User & Non & User & Non \\
\midrule
Age & 40.2 & 49.0 & 42.3 & 49.0 & 44.1 & 49.4 & 44.3 & 49.9 \\
Income (\$000) & 52.9 & 30.5 & 52.0 & 26.2 & 50.0 & 26.6 & 47.5 & 26.3 \\
College graduate & 0.43 & 0.21 & 0.45 & 0.21 & 0.42 & 0.20 & 0.37 & 0.20 \\
\bottomrule
\end{tabular}
\end{center}

\begin{itemize}
\tightlist
\item
  Internet adoption spread. The ``treated'' group got \textbf{older,
  poorer and less educated} over the window --- and older people buy
  less music. The composition of the treated group changed exactly when
  treatment arrived.
\item
  Nobody dropped out and nobody anticipated. Parallel trends breaks
  \textbf{because different people are sampled.}
\end{itemize}
\end{frame}

\begin{frame}{Diagnosing and repairing compositional change}
\protect\phantomsection\label{diagnosing-and-repairing-compositional-change}
\begin{itemize}
\tightlist
\item
  \textbf{Diagnostic 1.} A balance table \emph{over time}, as on the
  previous slide --- not just at baseline.
\item
  \textbf{Diagnostic 2.} Run the DiD with an exogenous covariate as the
  \emph{outcome}: \[
  X_{it} = \alpha + \beta_1 D_i + \beta_2 \text{Post}_t + \delta\,(D_i \times \text{Post}_t) + \varepsilon_{it}.
  \] Napster cannot make anyone older. A significant \(\hat\delta\) says
  the sample's composition shifted differentially at treatment. If that
  \(X\) also predicts \(Y^0\), parallel trends is mechanically broken.
\end{itemize}

\vspace{0.3em}

\begin{itemize}
\tightlist
\item
  \textbf{Repair (Hong 2013).} Estimate \textbf{two} propensity scores,
  one per period, \(P_b = \Pr(D=1 \mid T=0, X)\) and
  \(P_a = \Pr(D=1 \mid T=1, X)\); compute the four cell means with
  inverse-probability weights from the score of their own period; take
  the \(2 \times 2\) of the weighted means. Equivalently, the usual DiD
  regression by weighted least squares.
\item
  Panels are expensive; in much of the world the data are repeated
  cross-sections. This fix deserves more use than it gets (see also
  Sant'Anna and Xu 2023).
\end{itemize}
\end{frame}

\section{2. Triple Differences}\label{triple-differences}

\begin{frame}{Gruber (1994): maternity benefits and women's wages}
\protect\phantomsection\label{gruber-1994-maternity-benefits-and-womens-wages}
\begin{itemize}
\tightlist
\item
  Some states mandated that employer health insurance cover maternity.
  Who pays? If wages cannot adjust worker by worker, the mandate makes
  childbearing-age women more expensive to hire.
\item
  The obvious DiD: married women aged 20--40 in mandate states
  vs.~non-mandate states, before vs.~after. It identifies \[
  \hat\delta = \underbrace{D}_{\text{effect}} + \underbrace{(NJ_t - PA_t)}_{\text{non-parallel trends bias}},
  \] writing \(NJ\) for experimental states and \(PA\) for the rest, and
  \(NJ_t - PA_t\) for the difference in their untreated trends.
\item
  Gruber worried that state labour markets were on different paths for
  reasons unrelated to the mandate. Then the DiD is biased and should
  not be used as it stands.
\item
  His move: find a group in the \textbf{same states} that the law does
  not touch --- single men aged 20--40 and workers over 40 --- and run
  their DiD too.
\end{itemize}
\end{frame}

\begin{frame}{Triple differences is a design, not a falsification}
\protect\phantomsection\label{triple-differences-is-a-design-not-a-falsification}
The untouched group's DiD has no \(D\) in it:
\(\hat\delta_{\text{alt}} = (NJ_t - PA_t)\), so \[
\hat\delta_{DDD} = \hat\delta_{\text{primary}} - \hat\delta_{\text{alt}} = D + (NJ_t - PA_t) - (NJ_t - PA_t) = D.
\]

\begin{itemize}
\tightlist
\item
  The identifying assumption is \textbf{parallel bias} (Olden and Møen
  2022): the same non-parallel-trends bias for eligible and ineligible
  groups, \[
  \begin{aligned}
  &\mathbb{E}[\Delta Y^0 \mid D=1, G=1] - \mathbb{E}[\Delta Y^0 \mid D=0, G=1] \\
  &\qquad = \mathbb{E}[\Delta Y^0 \mid D=1, G=0] - \mathbb{E}[\Delta Y^0 \mid D=0, G=0],
  \end{aligned}
  \] with \(D\) the state and \(G\) the eligible group.
\item
  The alternative DiD is \textbf{not required to be zero}, only to equal
  the primary DiD's bias. If it \emph{were} zero, the primary DiD was
  fine all along, and the alternative group belonged in a ``same
  outcome, different group'' falsification, as with Miller, Johnson and
  Wherry's elderly.
\end{itemize}
\end{frame}

\begin{frame}{Triple differences by regression}
\protect\phantomsection\label{triple-differences-by-regression}
Eight averages and seven subtractions. Identical to OLS with
\textbf{all} interactions of period \(\tau_t\), group \(\delta_g\) and
state \(D_i\): \[
\begin{aligned}
Y_{igt} = \alpha &+ \beta_1\tau_t + \beta_2 \delta_g + \beta_3 D_i + \beta_4(\delta\tau)_{gt} \\
&+ \beta_5(\tau D)_{ti} + \beta_6(\delta D)_{ig} + \beta_7(\delta \tau D)_{igt} + \varepsilon_{igt}.
\end{aligned}
\]

\begin{itemize}
\tightlist
\item
  \(\beta_7\) is the triple difference. Drop any of the two-way terms
  and it no longer is.
\item
  Event-study version: interact the two-way and three-way terms with
  every period but \(\tau = -1\); the leads and lags of the three-way
  interaction are the triple-difference event study.
\item
  Cluster where the law varies: the state. Show event studies of
  \textbf{both} DiDs and of their difference; the two pre-trends should
  look alike. Almost nobody shows them; show yours.
\end{itemize}
\end{frame}

\begin{frame}{Gruber's table, and what an event study adds}
\protect\phantomsection\label{grubers-table-and-what-an-event-study-adds}
\begin{center}\scriptsize
\begin{tabular}{lccc}
\toprule
Log hourly wage & Before & After & Difference \\
\midrule
\multicolumn{4}{l}{\emph{A. Married women 20--40}} \\
Experimental states & 1.547 & 1.513 & $-0.034$ \\
Non-experimental states & 1.369 & 1.397 & $0.028$ \\
Difference-in-differences & & & $\mathbf{-0.062}$ (0.022) \\
\midrule
\multicolumn{4}{l}{\emph{B. Over 40 and single men 20--40}} \\
Experimental states & 1.759 & 1.748 & $-0.011$ \\
Non-experimental states & 1.630 & 1.627 & $-0.003$ \\
Difference-in-differences & & & $\mathbf{-0.008}$ (0.014) \\
\midrule
Triple difference & & & $\mathbf{-0.054}$ (0.026) \\
\bottomrule
\end{tabular}
\end{center}

\begin{itemize}
\tightlist
\item
  Panel B \emph{estimates the bias} of panel A. Here it is small, so DDD
  and DiD nearly agree: married women's wages fell about \(5\) percent.
  The mandate was paid for by the workers it protected.
\item
  No event study in Gruber (1994); the practice did not arrive until the
  2000s.
\end{itemize}
\end{frame}

\section{3. Covariates}\label{covariates}

\begin{frame}{Levels, trends, and why covariates enter}
\protect\phantomsection\label{levels-trends-and-why-covariates-enter}
\begin{itemize}
\tightlist
\item
  Parallel trends is about trends, not levels. Two groups at different
  levels are fine \textbf{in principle}.
\item
  In practice a large level gap is a warning. Something made the groups
  different; could the same thing make their untreated \emph{trends}
  different? Kahn-Lang and Lang (2019): every DiD paper should say why
  the levels differ and why that would not affect the trends.
\item
  Often it is the \textbf{distribution of covariates} that generates
  both. Then condition on them: \[
  \mathbb{E}[Y^0_{\text{post}} - Y^0_{\text{pre}} \mid D=1, X] \;=\; \mathbb{E}[Y^0_{\text{post}} - Y^0_{\text{pre}} \mid D=0, X]
  \] --- \textbf{conditional parallel trends} (Heckman, Ichimura and
  Todd 1997). Parallel trends among units with the same \(X\).
\item
  Two things come with it: \textbf{common support} --- for every treated
  \(X\) there must be controls with that \(X\) --- or a functional form
  that extrapolates where there are none; and a decision about
  \textbf{which} \(X\).
\end{itemize}
\end{frame}

\begin{frame}{Which covariates}
\protect\phantomsection\label{which-covariates}
\begin{itemize}
\tightlist
\item
  You are imputing \(\mathbb{E}[Y^0_{\text{post}} \mid D=1]\) from
  controls \emph{with the same \(X\)}. So \(X\) must do two things:

  \begin{enumerate}
  \tightlist
  \item
    \(X \rightarrow D\): it is distributed differently across treated
    and controls;
  \item
    \(X \rightarrow \Delta \mathbb{E}[Y^0]\): it drives the untreated
    \textbf{trend}.
  \end{enumerate}
\item
  A covariate that does only the first is harmless but useless; one that
  does only the second adds precision but is not needed for
  identification. Both at once is the confounder of trends.
\item
  Where the list comes from: the field that studies ordinary trends in
  your outcome; a DAG that forces you to say what moves \(Y^0\); or a
  data-driven selection fitted \textbf{on untreated units only}
  (Borgschulte and Vogler 2020 use LASSO this way).
\item
  \textbf{Never} a post-treatment covariate: it may itself respond to
  treatment. Measure \(X\) at baseline.
\item
  The list can swing the estimate a lot. Fix it before you see results,
  and say so.
\end{itemize}
\end{frame}

\begin{frame}{Is there imbalance? The normalized difference}
\protect\phantomsection\label{is-there-imbalance-the-normalized-difference}
\[
\text{Norm.\ Diff}_\omega = \frac{\bar X_{\omega,T} - \bar X_{\omega,C}}{\sqrt{(S^2_{\omega,T} + S^2_{\omega,C})/2}}
\]

Unit-free; rule of thumb \(|{\cdot}| < 0.25\) (Imbens and Rubin 2015).
Lecture 3's LaLonde data, treated vs.~CPS controls:

\begin{center}\small
\begin{tabular}{lccc}
\toprule
 & Mean treated & Mean control & Norm.\ diff. \\
\midrule
Age & 25.8 & 33.2 & $-0.80$ \\
Education & 10.3 & 12.0 & $-0.68$ \\
Black & 0.84 & 0.07 & $2.43$ \\
Unemployed 1974 & 0.71 & 0.12 & $1.49$ \\
\bottomrule
\end{tabular}
\end{center}

\begin{itemize}
\tightlist
\item
  Everything but Hispanic is far past the line. Irrelevant if these
  variables do not drive trends in earnings; fatal if they do. Either
  way, unconditional parallel trends is not a claim you could make with
  a straight face: trade it for the conditional version and an estimator
  that uses it.
\end{itemize}
\end{frame}

\begin{frame}{Estimator 1: inverse probability weighting (Abadie 2005)}
\protect\phantomsection\label{estimator-1-inverse-probability-weighting-abadie-2005}
\begin{itemize}
\tightlist
\item
  Estimate \(p(X) = \Pr(D=1 \mid X_b)\) from \textbf{baseline}
  covariates. Then \[
  \hat\delta_{IPW} = \mathbb{E}\!\left[\frac{Y_{\text{post}} - Y_{\text{pre}}}{\Pr(D=1)} \cdot \frac{D - p(X)}{1 - p(X)}\right].
  \]
\item
  The weight is \(1\) for every treated unit and \(-p/(1-p)\) for a
  control. \textbf{Only the control group is reweighted}, to look like
  the treated group in \(X\); for the ATT that is all that is needed.
\end{itemize}

\begin{center}\scriptsize
\begin{tabular}{lcccrr}
\toprule
 & $\Delta Y$ & $D$ & $p(X)$ & weight & weighted $\Delta Y$ \\
\midrule
Pat & 45 & 0 & 0.75 & $-3$ & $-135$ \\
Tiffany & 100 & 1 & 0.90 & $1$ & $100$ \\
Dolores & 42 & 0 & 0.25 & $-1/3$ & $-14$ \\
Danny & 75 & 1 & 0.50 & $1$ & $75$ \\
\bottomrule
\end{tabular}
\end{center}

\begin{itemize}
\tightlist
\item
  Assumptions: no anticipation, conditional parallel trends, common
  support, a \textbf{correctly specified score}.
\item
  Outliers here are \textbf{extreme weights}: a control with
  \(p = 0.99991\) gets weight \(-11{,}110\). Count them before you
  estimate; trim if you must.
\end{itemize}
\end{frame}

\begin{frame}{Estimator 2: outcome regression (Heckman, Ichimura and
Todd 1997)}
\protect\phantomsection\label{estimator-2-outcome-regression-heckman-ichimura-and-todd-1997}
\begin{itemize}
\tightlist
\item
  The mirror image. Using \textbf{controls only}, fit the change in
  outcomes on baseline covariates: \(\hat\mu(X)\) predicts
  \(\Delta Y^0\) from \(X\).
\item
  Apply the fitted function to each treated unit's \(X\) to impute its
  counterfactual change. Then \[
  \hat\delta_{OR} = \big(\bar Y_{1,\text{post}} - \bar Y_{1,\text{pre}}\big) - \frac{1}{n_T}\sum_{i: D_i = 1}\big[\hat\mu_{0,\text{post}}(X_i) - \hat\mu_{0,\text{pre}}(X_i)\big].
  \]
\item
  Regression adjustment, as in Lecture 3, applied to the long
  difference. Under conditional parallel trends the control group's
  \(\Delta Y^0 \sim X\) relationship is the treated group's too; that is
  the entire content of the assumption.
\item
  Assumptions: the same three, plus a \textbf{correctly specified
  outcome model} in place of a correctly specified score.
\item
  Notice the asymmetry. IPW touches the \emph{controls} and leaves the
  treated alone; OR touches the \emph{treated} and leaves the controls
  alone. Different transformations, same target.
\end{itemize}
\end{frame}

\begin{frame}{Estimator 3: doubly robust (Sant'Anna and Zhao 2020)}
\protect\phantomsection\label{estimator-3-doubly-robust-santanna-and-zhao-2020}
\begin{itemize}
\tightlist
\item
  Do both. Reweight the controls with the score \textbf{and} subtract
  the outcome-model prediction from everyone: \[
  \hat\delta_{DR} = \mathbb{E}\!\left[\left(\frac{D}{\mathbb{E}[D]} - \frac{\frac{p(X)(1-D)}{1-p(X)}}{\mathbb{E}\!\left[\frac{p(X)(1-D)}{1-p(X)}\right]}\right)\big(\Delta Y - \hat\mu_{0,\Delta}(X)\big)\right].
  \]
\item
  Consistent if \textbf{either} the score or the outcome model is right.
  You are allowed to be wrong --- once.
\item
  The cost is nothing you were not already paying: no anticipation,
  conditional parallel trends, common support, and one of two models.
\end{itemize}

\vspace{0.3em}

\begin{itemize}
\tightlist
\item
  Lecture 3b derived why the repair term cancels. Here the only change
  is that the outcome is a \emph{change}, and the assumption is about
  \emph{trends}.
\item
  Chang (2020) does the same with machine-learned nuisance functions;
  worth knowing when the covariate list is long.
\end{itemize}
\end{frame}

\begin{frame}{Estimator 4: the regression everyone runs}
\protect\phantomsection\label{estimator-4-the-regression-everyone-runs}
\[
Y_{it} = \alpha_1 + \alpha_2\,\text{Post}_t + \alpha_3 D_i + \delta\,(D_i \times \text{Post}_t) + \theta X_{it} + \varepsilon_{it}
\]

\begin{itemize}
\tightlist
\item
  Nothing stops you. But two more assumptions are hiding in
  \(\theta X_{it}\).
\item
  \textbf{Homogeneous effects.} Let \(X\) affect \(Y^1\) by \(\theta_1\)
  and \(Y^0\) by \(\theta_2\). Then \[
  \text{ATT} = \mathbb{E}[Y^1_{\text{post}} \mid D=1, X] - \mathbb{E}[Y^0_{\text{post}} \mid D=1, X] = \delta + (\theta_1 - \theta_2)X .
  \] The regression reports \(\delta\), which is the ATT only if
  \(\theta_1 = \theta_2\): the same effect for old and young, rich and
  poor. Is the age--earnings profile the same with and without a degree?
  That is the assumption.
\item
  \textbf{No covariate-specific trends.} With cell means
  \(\bar X_{dt}\), the DiD of the cells is \[
  \delta^{DD} = \delta + \theta(\bar X_{11} - \bar X_{10}) - \theta(\bar X_{01} - \bar X_{00}),
  \] and the regression imposes one \(\theta\) in every period. If the
  \emph{return} to \(X\) changes over time, or \(X\) trends differently
  across groups, the extra terms do not cancel.
\item
  Plus functional form, since nothing requires common support to stop
  extrapolation.
\end{itemize}
\end{frame}

\begin{frame}{The four estimators, simulated}
\protect\phantomsection\label{the-four-estimators-simulated}
\begin{center}
\includegraphics[width=0.78\linewidth,height=\textheight,keepaspectratio]{07b_did_complex_designs_files/figure-beamer/cov-sim-1.pdf}
\end{center}

\begin{itemize}
\tightlist
\item
  Two periods, \(1{,}000\) units, \(400\) replications. Treatment
  depends on age and GPA; so does the untreated trend, with the
  \emph{return} to each covariate doubling over time; and the effect
  varies with both. Conditional parallel trends holds and both nuisance
  models are correct.
\item
  Mean bias: TWFE \(807\); IPW \(5\); OR \(-5\); DR \(-7\). Standard
  deviations \(144\), \(242\), \(159\), \(165\): IPW pays in variance
  for near-one scores among controls.
\item
  The problem is not regression but the \textbf{additive, time-constant}
  way this specification lets \(X\) in. Wooldridge (2021) fixes it with
  OLS; the three estimators above fix it by construction.
\end{itemize}
\end{frame}

\section{4. Differential Timing and Two-Way Fixed
Effects}\label{differential-timing-and-two-way-fixed-effects}

\begin{frame}{When ``Post'' has no meaning}
\protect\phantomsection\label{when-post-has-no-meaning}
Card and Krueger again, but suppose New York had also raised its minimum
wage, in 1993.

\begin{center}\scriptsize
\begin{tabular}{lccc|lccc|lccc}
\toprule
Year & State & Treat & Post & Year & State & Treat & Post & Year & State & Treat & Post \\
\midrule
1991 & NJ & 1 & 0 & 1991 & NY & 1 & 0 & 1991 & PA & 0 & 0 \\
1992 & NJ & 1 & 1 & 1992 & NY & 1 & 0 & 1992 & PA & 0 & ? \\
1993 & NJ & 1 & 1 & 1993 & NY & 1 & 1 & 1993 & PA & 0 & ? \\
\bottomrule
\end{tabular}
\end{center}

\begin{itemize}
\tightlist
\item
  Post relative to \emph{which} date? The interaction regression cannot
  even be written down. So everyone runs the version that can: \[
  Y_{ist} = \beta_0 + \delta D_{ist} + \tau_t + \sigma_s + \varepsilon_{ist}, \qquad D_{ist} = \mathbf{1}\{\text{unit } s \text{ is treated at } t\}.
  \]
\item
  Coding \(D\) is trivial. The question is what \(\hat\delta^{TWFE}\)
  \emph{is}, in potential outcomes.
\end{itemize}
\end{frame}

\begin{frame}{Every \(2 \times 2\) at once (Goodman-Bacon 2021)}
\protect\phantomsection\label{every-2-times-2-at-once-goodman-bacon-2021}
\begin{itemize}
\tightlist
\item
  By Frisch--Waugh--Lovell, \(\hat\delta^{TWFE}\) is the slope of \(Y\)
  on \(D\) after both have been demeaned by unit and by period:
  \(\hat\delta = \sum \tilde D_{st}\tilde Y_{st} \big/ \sum \tilde D_{st}^2\).
\item
  Goodman-Bacon showed that this slope is \textbf{exactly} a weighted
  average of all the simple \(2 \times 2\) DiDs the data contain. Three
  timing groups \(a < b < c\) and a never-treated \(U\) give nine:
\end{itemize}

\begin{center}\small
\begin{tabular}{ccc}
\toprule
$a$ vs $b$ & $b$ vs $a$ & $c$ vs $a$ \\
$a$ vs $c$ & $b$ vs $c$ & $c$ vs $b$ \\
$a$ vs $U$ & $b$ vs $U$ & $c$ vs $U$ \\
\bottomrule
\end{tabular}
\end{center}

\begin{itemize}
\tightlist
\item
  ``\(a\) vs \(b\)'' uses \(b\) as a control while it is still
  untreated: a \textbf{not-yet-treated} \(2 \times 2\). Fine.
\item
  ``\(b\) vs \(a\)'' uses \(a\) as a control \textbf{after \(a\) was
  treated}. That is the already-treated comparison of Lecture 7a --- the
  one that subtracts the control's own evolving effect. Three of the
  nine are of this kind. Nobody chose them; OLS did.
\end{itemize}
\end{frame}

\begin{frame}{The three kinds of \(2 \times 2\)}
\protect\phantomsection\label{the-three-kinds-of-2-times-2}
Two timing groups, \(k\) before \(l\), and a never-treated \(U\);
\emph{mid} is the window between the two dates: \[
\begin{aligned}
\hat\delta^{2\times2}_{kU} &= (\bar y_k^{\,post(k)} - \bar y_k^{\,pre(k)}) - (\bar y_U^{\,post(k)} - \bar y_U^{\,pre(k)}) &&\text{never-treated control: legal} \\[2pt]
\hat\delta^{2\times2}_{kl} &= (\bar y_k^{\,mid} - \bar y_k^{\,pre(k)}) - (\bar y_l^{\,mid} - \bar y_l^{\,pre(k)}) &&\text{not-yet-treated control: legal} \\[2pt]
\hat\delta^{2\times2}_{lk} &= (\bar y_l^{\,post(l)} - \bar y_l^{\,mid}) - (\bar y_k^{\,post(l)} - \bar y_k^{\,mid}) &&\text{already-treated control: forbidden}
\end{aligned}
\]

\begin{itemize}
\tightlist
\item
  In the third, \(k\)'s outcomes in both windows are \(Y^1\). Substitute
  potential outcomes and, as in 7a, \[
  \hat\delta^{2\times2}_{lk} = \text{ATT}_l(post) + \text{PT bias} - \big[\text{ATT}_k(post(l)) - \text{ATT}_k(mid)\big].
  \] If \(k\)'s effect was still growing, the growth is subtracted from
  \(l\)'s estimate.
\item
  ``Forbidden comparison'' is Borusyak, Jaravel and Spiess's phrase.
  Independently found by de Chaisemartin and D'Haultfoeuille (2020),
  Goodman-Bacon (2021) and Sun and Abraham (2021), all early-career
  authors, all within a few years.
\end{itemize}
\end{frame}

\begin{frame}{The weights: group size times treatment variance}
\protect\phantomsection\label{the-weights-group-size-times-treatment-variance}
\[
\hat\delta^{TWFE} = \sum_{k \ne U} s_{kU}\,\hat\delta^{2\times2}_{kU} + \sum_{k \ne U}\sum_{l > k}\big[s^{k}_{kl}\,\hat\delta^{2\times2,k}_{kl} + s^{l}_{kl}\,\hat\delta^{2\times2,l}_{lk}\big], \qquad \sum s = 1,
\] \[
s_{kU} = \frac{(n_k + n_U)^2\; n_{kU}(1 - n_{kU})\; \bar D_k(1 - \bar D_k)}{\widehat{\text{Var}}(\tilde D)}, \quad n_{kU} = \frac{n_k}{n_k + n_U},
\] and similarly for the timing pairs. Every weight is
\textbf{positive}. Two ingredients:

\begin{itemize}
\tightlist
\item
  \textbf{Sample shares} \(n_k\): bigger groups count more, as you would
  expect.
\item
  \textbf{Treatment variance} \(\bar D_k(1-\bar D_k)\), where
  \(\bar D_k\) is the share of the panel a group spends treated. It
  peaks at \(\bar D = 0.5\): \(0.1 \times 0.9 = 0.09\),
  \(0.4 \times 0.6 = 0.24\), \(0.5 \times 0.5 = 0.25\).
\item
  So TWFE \textbf{over-weights groups treated near the middle of the
  panel} and under-weights early and late adopters --- not because their
  comparisons are better, but because OLS likes variance. Change the
  panel's start or end date and the weights change with it.
\end{itemize}
\end{frame}

\begin{frame}{From arithmetic to causality}
\protect\phantomsection\label{from-arithmetic-to-causality}
Replace each \(2 \times 2\) by its potential-outcome content, with no
anticipation: \[
\operatorname{plim}\hat\delta^{TWFE} \;=\; \text{VWATT} \;+\; \text{VWPT} \;-\; \Delta\text{ATT}.
\]

\begin{itemize}
\tightlist
\item
  \textbf{VWATT}: a variance-weighted average of the group ATTs. Causal,
  but weighted by the panel's start and end dates rather than by
  anything you chose.
\item
  \textbf{VWPT}: variance-weighted parallel-trends bias. Zero if
  parallel trends holds for every pair.
\item
  \(\boldsymbol{\Delta}\)\textbf{ATT}: the already-treated controls'
  \emph{change} in effect, entering with a \textbf{minus sign}. Zero
  only if effects are constant over time.
\end{itemize}

\vspace{0.3em}

\begin{itemize}
\tightlist
\item
  Parallel trends is not enough for this regression. You also need the
  effect \textbf{not to change over time.} If it grows, the estimate is
  biased toward zero; if it grows enough, the sign flips.
  \(\text{VWATT} = 10\) and \(\Delta\text{ATT} = 3\) gives \(7\);
  \(\Delta\text{ATT} = 11\) gives \(-1\). Rat poison for a headache,
  because the regression said so.
\end{itemize}
\end{frame}

\begin{frame}{The Baker simulation: the truth}
\protect\phantomsection\label{the-baker-simulation-the-truth}
\begin{center}
\includegraphics[width=0.76\linewidth,height=\textheight,keepaspectratio]{07b_did_complex_designs_files/figure-beamer/baker-raw-1.pdf}
\end{center}

\begin{itemize}
\tightlist
\item
  \(1{,}000\) firms, \(1980\)--\(2009\), four equal cohorts treated in
  \(1986\), \(1992\), \(1998\), \(2004\). Every firm's untreated outcome
  grows by \(1\) a year plus small noise: parallel trends holds
  \textbf{firm by firm}. No anticipation, no omitted variables,
  homoskedastic errors.
\item
  Two outcomes. \textbf{Constant} effects: \(10\), \(8\), \(6\), \(4\)
  by cohort, forever. \textbf{Dynamic} effects: the same numbers,
  growing by that amount every year treated --- the 1986 cohort is
  \(240\) above its counterfactual by 2009. The figure shows the dynamic
  one.
\end{itemize}
\end{frame}

\begin{frame}{The Baker simulation: the parameters and the regression}
\protect\phantomsection\label{the-baker-simulation-the-parameters-and-the-regression}
\begin{center}\scriptsize
\begin{tabular}{lcccccc|c}
\toprule
 & ATT(1986) & ATT(1992) & ATT(1998) & ATT(2004) & Simple & Group & $\hat\delta^{TWFE}$ (s.e.) \\
\midrule
Constant & 10 & 8 & 6 & 4 & 8 & 7 & $7.00$ ($0.027$) \\
Dynamic & 125 & 76 & 39 & 14 & 82 & 63.5 & $\mathbf{-6.69}$ ($0.670$) \\
\bottomrule
\end{tabular}
\end{center}

\begin{itemize}
\tightlist
\item
  \emph{Simple ATT}: average over all treated group-years. \emph{Group
  ATT}: average of the four cohort averages. Both are correct; they
  weight the cells differently.
\item
  \textbf{Constant effects:} \(\Delta\text{ATT} = 0\) and parallel
  trends holds, so only the variance weights remain. Middle cohorts get
  more weight than the edges, and the answer is \(7\) --- a VWATT with a
  causal reading, though not the simple ATT.
\item
  \textbf{Dynamic effects:} every effect positive, every weight
  positive, parallel trends exact, errors clean --- and the regression
  returns \(\mathbf{-6.69}\) with a \(t\)-statistic of \(10\).
  \textbf{The sign is wrong.} Strict exogeneity holds and the answer is
  still wrong.
\end{itemize}
\end{frame}

\begin{frame}{The Baker simulation: the decomposition}
\protect\phantomsection\label{the-baker-simulation-the-decomposition}
Bacon's decomposition, computed by hand from group means --- no package
--- reproduces the coefficient exactly:

\begin{center}\small
\begin{tabular}{lcc}
\toprule
 & Weight & Average $2 \times 2$ \\
\midrule
Earlier treated vs later control (not-yet-treated) & $0.50$ & $51.8$ \\
Later treated vs earlier control (already-treated) & $0.50$ & $-65.2$ \\
\midrule
$0.50 \times 51.8 + 0.50 \times (-65.2)$ & & $-6.69$ \\
\bottomrule
\end{tabular}
\end{center}

\begin{itemize}
\tightlist
\item
  No never-treated group, so every \(2 \times 2\) is between two
  cohorts: \(12\) of them. Half the weight sits on comparisons whose
  control was \textbf{already treated and still climbing}; those average
  \(-65.2\).
\item
  Under constant effects the two halves give \(8.4\) and \(5.6\):
  harmless, because a control whose effect does not move differences out
  like a fixed effect.
\item
  The weights are \textbf{observable}. Compute them for your design
  before deciding whether TWFE is good enough.
\end{itemize}
\end{frame}

\begin{frame}[fragile]{Bacon on real data: castle doctrine, all cohorts}
\protect\phantomsection\label{bacon-on-real-data-castle-doctrine-all-cohorts}
\begin{center}
\includegraphics[width=0.76\linewidth,height=\textheight,keepaspectratio]{07b_did_complex_designs_files/figure-beamer/castle-bacon-1.pdf}
\end{center}

\begin{itemize}
\tightlist
\item
  Lecture 7a used the 2006 cohort against the never-treated states. Now
  all \(21\) adopting states, \(2005\)--\(2009\), plus \(29\) never
  adopters. TWFE on \(D_{st} = \mathbf{1}\{t \ge \text{adoption}\}\)
  gives \(\hat\delta = 0.082\) (\(0.059\)).
\item
  Inside it: \(25\) simple \(2 \times 2\)'s. Weight on treated-vs-never:
  \(0.91\) (average \(0.088\)); on earlier-vs-later: \(0.06\)
  (\(-0.006\)); on the forbidden later-vs-earlier: \(0.03\) (\(0.070\)).
  The 2006-vs-never cell alone carries \(0.59\).
\item
  With a large never-treated group the forbidden comparisons get little
  weight and TWFE is close to the clean answer. That is \textbf{luck of
  the design}, not a property of the estimator; check it every time.
  Stata: \texttt{bacondecomp}; R: \texttt{bacondecomp}.
\end{itemize}
\end{frame}

\section{5. Estimators That Refuse the Forbidden
Comparison}\label{estimators-that-refuse-the-forbidden-comparison}

\begin{frame}{Building blocks: the group-time ATT (Callaway and
Sant'Anna 2021)}
\protect\phantomsection\label{building-blocks-the-group-time-att-callaway-and-santanna-2021}
\begin{itemize}
\tightlist
\item
  Stop asking for one number. For the cohort first treated at \(g\) and
  calendar period \(t\),
  \(\text{ATT}(g, t) = \mathbb{E}[Y^1_t - Y^0_t \mid G_g = 1]\).
\item
  Estimate each one as its own \(2 \times 2\): cohort \(g\) from its
  baseline \(g-1\) to \(t\), against units untreated in both periods ---
  \textbf{never-treated}, or \textbf{not-yet-treated}. Never a unit that
  is already treated.
\item
  With covariates, the \(2 \times 2\) becomes the IPW, OR or DR
  estimator of Section 3, with a score
  \(\hat p(X) = \Pr(G_g = 1 \mid X, G_g + C = 1)\) estimated among
  cohort \(g\) and its controls \(C\): \[
  \text{ATT}(g,t) = \mathbb{E}\!\left[\left(\frac{G_g}{\mathbb{E}[G_g]} - \frac{\frac{\hat p(X) C}{1 - \hat p(X)}}{\mathbb{E}\!\left[\frac{\hat p(X) C}{1 - \hat p(X)}\right]}\right)\big(Y_t - Y_{g-1}\big)\right].
  \]
\item
  Assumptions: \textbf{absorbing} treatment; no anticipation; parallel
  trends given \(X\) for the chosen controls; weak overlap; one correct
  nuisance model if \(X\) is used. Then \textbf{aggregate} with weights
  you choose and can defend.
\end{itemize}
\end{frame}

\begin{frame}{Which cells exist}
\protect\phantomsection\label{which-cells-exist}
Baker data, dynamic effects. A cell \(\text{ATT}(g,t)\) is estimable
only while some unit is still untreated at \(t\):

\begin{center}\scriptsize
\begin{tabular}{lccccc}
\toprule
Year & $g = 1986$ & $g = 1992$ & $g = 1998$ & $g = 2004$ & Untreated control? \\
\midrule
1985 & 0 & 0 & 0 & 0 & yes \\
1986 & 10 & 0 & 0 & 0 & yes \\
1992 & 70 & 8 & 0 & 0 & yes \\
1998 & 130 & 56 & 6 & 0 & yes \\
2003 & 180 & 96 & 36 & 0 & yes \\
2004 & 190 & 104 & 42 & 4 & \textbf{no} \\
2009 & 240 & 144 & 72 & 24 & \textbf{no} \\
\midrule
ATT($g$), all cells & 125 & 76 & 39 & 14 & \\
ATT($g$), feasible cells & 95 & 52 & 21 & -- & \\
\bottomrule
\end{tabular}
\end{center}

\begin{itemize}
\tightlist
\item
  After 2003 everyone is treated; the 2004 cohort exists only as a
  control. The feasible simple ATT is \(68.3\) and the feasible group
  ATT \(56\), below \(82\) and \(63.5\): the late, large effects are not
  identified by \textbf{any} honest comparison. TWFE used those years
  anyway, through forbidden comparisons.
\end{itemize}
\end{frame}

\begin{frame}[fragile]{One cell by hand, then all of them}
\protect\phantomsection\label{one-cell-by-hand-then-all-of-them}
\(\text{ATT}(1986, 1986)\): the 1986 cohort's mean outcome in 1986 minus
in 1985, less the same change for firms not yet treated in 1986 --- the
1992, 1998 and 2004 cohorts: \[
\widehat{\text{ATT}}(1986,1986) = (\bar y_{1986}^{\,g=1986} - \bar y_{1985}^{\,g=1986}) - (\bar y_{1986}^{\,C} - \bar y_{1985}^{\,C}) = 10.03 \qquad (\text{truth } 10).
\]

\begin{itemize}
\tightlist
\item
  Four averages and three subtractions. With covariates they become four
  \emph{weighted} averages. That is the whole estimator; the package
  only loops.
\item
  All \(36\) feasible post-treatment cells, aggregated:
\end{itemize}

\begin{center}\small
\begin{tabular}{lccc}
\toprule
 & Estimate & Feasible truth & Full truth \\
\midrule
Simple ATT: mean over cells & $68.34$ & $68.33$ & $82$ \\
Group ATT: mean of ATT($g$) & $56.00$ & $56$ & $63.5$ \\
\bottomrule
\end{tabular}
\end{center}

\begin{itemize}
\tightlist
\item
  ATT(\(g\)) by cohort: \(95.0\), \(52.0\), \(21.0\). Stata:
  \texttt{csdid2\ y,\ ivar(id)\ time(year)\ gvar(treat\_date)\ notyet\ method(stdipw)\ long2};
  R: \texttt{did::att\_gt}.
\end{itemize}
\end{frame}

\begin{frame}{Event studies from building blocks}
\protect\phantomsection\label{event-studies-from-building-blocks}
\begin{center}
\includegraphics[width=0.76\linewidth,height=\textheight,keepaspectratio]{07b_did_complex_designs_files/figure-beamer/baker-es-1.pdf}
\end{center}

\begin{itemize}
\tightlist
\item
  Average the cells by \textbf{relative time} \(\ell = t - g\). At
  \(\ell = 0\) the three estimable cohorts give \((10 + 8 + 6)/3 = 8\);
  at \(\ell = 6\) only 1986 and 1992 remain, \((70 + 56)/2 = 63\); from
  \(\ell = 12\) only 1986 is left. The staircase is real: \textbf{later
  lags are early adopters' effects}, not the population's.
\item
  Same sample for all three: years to 2003, with the 2004 cohort as the
  untreated control. CS and SA sit on the truth. The TWFE event study
  shows a \emph{declining} pre-trend and near-zero post effects in data
  where neither exists. The next slides explain why.
\end{itemize}
\end{frame}

\begin{frame}{Castle doctrine, all cohorts: CS against TWFE}
\protect\phantomsection\label{castle-doctrine-all-cohorts-cs-against-twfe}
\begin{center}\small
\begin{tabular}{lcccccccc}
\toprule
 & Static & $\ell=-3$ & $\ell=-2$ & $\ell=0$ & $\ell=1$ & $\ell=2$ & $\ell=3$ & $\ell=4$ \\
\midrule
TWFE & $0.082$ & $0.054$ & $0.059$ & $0.092$ & $0.106$ & $0.115$ & $0.110$ & $0.084$ \\
CS, never-treated & $0.110$ & $0.057$ & $0.058$ & $0.097$ & $0.112$ & $0.112$ & $0.137$ & $0.093$ \\
\bottomrule
\end{tabular}
\end{center}

\begin{itemize}
\tightlist
\item
  CS: \(20\) group-time cells against the \(29\) never-adopting states,
  aggregated with cohort sizes as weights (\(1, 13, 4, 2, 1\) states).
  Event-time cells use the fixed baseline \(g-1\).
\item
  The static CS estimate is \(11\) percent against TWFE's \(8\). The
  Bacon weights told us why: \(91\) percent of TWFE is clean
  comparisons, and the rest average lower. The event-study coefficients
  barely move.
\item
  Here the two agree on the story. In Baker, Larcker and Wang's
  replications they did not. You cannot know which case you are in
  without looking.
\end{itemize}
\end{frame}

\begin{frame}{Event studies with TWFE: what Sun and Abraham found}
\protect\phantomsection\label{event-studies-with-twfe-what-sun-and-abraham-found}
The canonical event study under differential timing,
\(Y_{it} = \alpha_i + \lambda_t + \sum_{\ell \ne -1}\mu_\ell\,\mathbf{1}\{t - E_i = \ell\} + \varepsilon_{it}\).

\begin{itemize}
\tightlist
\item
  Three assumptions: \textbf{parallel trends} for every cohort;
  \textbf{no anticipation}; a \textbf{homogeneous effect profile} ---
  the same path of effects after treatment for every cohort. Dynamics
  are fine, so long as everyone has the \emph{same} dynamics.
\item
  Sun and Abraham (2021) decompose each \(\hat\mu_\ell\) into
  cohort-specific effects \(\text{CATT}_{e,\ell}\):
\end{itemize}

\[
\hat\mu_\ell = \underbrace{\sum_e w^{\ell}_{e,\ell}\,\text{CATT}_{e,\ell}}_{\text{own period: weights sum to } 1} \;+\;
\underbrace{\sum_{\ell' \ne \ell}\sum_e w^{\ell}_{e,\ell'}\,\text{CATT}_{e,\ell'}}_{\text{other included periods: sum to } 0} \;+\;
\underbrace{\sum_{\ell' \in \text{excl}}\sum_e w^{\ell}_{e,\ell'}\,\text{CATT}_{e,\ell'}}_{\text{excluded periods: sum to } -1}.
\]

\begin{itemize}
\tightlist
\item
  Weights that sum to zero contain negatives: the coefficient for
  \(\ell = -2\) contains effects from \(\ell = +5\).
  \textbf{Contamination}, as with judge dummies in Lecture 5b
  (Goldsmith-Pinkham, Hull and Kolesár 2024).
\end{itemize}
\end{frame}

\begin{frame}{The two-cohort example}
\protect\phantomsection\label{the-two-cohort-example}
Two periods, two cohorts: cohort 1 treated in period 1, cohort 2 in
period 2. Include dummies for \(\ell \in \{-2, 0\}\), exclude
\(\{-1, +1\}\). Then \[
\hat\mu_{-2} = \underbrace{\text{CATT}_{2,-2}}_{\text{own}} + \underbrace{\tfrac12\text{CATT}_{1,0} - \tfrac12\text{CATT}_{2,0}}_{\text{other included}} + \underbrace{\tfrac12\text{CATT}_{1,1} - \text{CATT}_{1,-1} - \tfrac12\text{CATT}_{2,-1}}_{\text{excluded}}.
\]

\begin{itemize}
\tightlist
\item
  \textbf{Parallel trends} makes every term a causal effect rather than
  a trend difference.
\item
  \textbf{No anticipation} deletes the three pre-period terms.
\item
  \textbf{Homogeneous profiles} make
  \(\text{CATT}_{1,0} = \text{CATT}_{2,0}\), so the middle cancels.
\item
  What survives is \(\tfrac12\text{CATT}_{1,1}\): a ``pre-trend''
  coefficient made of a post-treatment effect, from a period you
  excluded. Nobody excludes a post period in practice; the point is the
  \emph{mechanism}. With heterogeneous profiles the middle term survives
  too, and that is what drew the sloping pre-trend in the Baker figure.
\end{itemize}
\end{frame}

\begin{frame}[fragile]{The interaction-weighted estimator}
\protect\phantomsection\label{the-interaction-weighted-estimator}
\begin{enumerate}
\tightlist
\item
  Run a \(2 \times 2\) (with or without covariates) for \textbf{each
  cohort and each relative period}, against the \textbf{never-treated}
  or the \textbf{last-treated} cohort. Not the not-yet-treated: that is
  the one difference from CS. These are the \(\text{CATT}_{e,\ell}\).
\item
  Weight each cohort by its share of the units observed at relative time
  \(\ell\).
\item
  Sum. The result is a properly weighted average of cohort effects at
  \(\ell\), and consistent for \(\text{ATT}_\ell\) under parallel trends
  and no anticipation alone.
\end{enumerate}

\vspace{0.3em}

\begin{itemize}
\tightlist
\item
  In the Baker figure this gives the truth to the first decimal, with
  the 2004 cohort as control. R:
  \texttt{fixest::sunab(cohort,\ period)}; Stata:
  \texttt{eventstudyinteract}.
\item
  CS and SA are the same building blocks with different default controls
  and aggregations. They can disagree when parallel trends holds for one
  control group and not the other --- which is a statement about your
  design, not about the estimators.
\end{itemize}
\end{frame}

\begin{frame}[fragile]{Treatments that switch off: de Chaisemartin and
D'Haultfoeuille (2020)}
\protect\phantomsection\label{treatments-that-switch-off-de-chaisemartin-and-dhaultfoeuille-2020}
\begin{itemize}
\tightlist
\item
  dCDH diagnose the same disease at the level of \textbf{unit-time
  effects}: TWFE is a weighted sum of every treated cell's effect, and
  some weights are \textbf{negative}, most often in periods when many
  units are treated and for units treated for long. If you must report
  TWFE, report the share of negative weights (Stata:
  \texttt{twowayfeweights}).
\item
  Their estimator, \(\text{DID}_M\), averages two kinds of
  \(2 \times 2\) at each date:

  \begin{itemize}
  \tightlist
  \item
    \textbf{joiners}: units switching into treatment at \(t\), against
    units untreated at \(t-1\) and \(t\);
  \item
    \textbf{leavers}: units switching out, against units treated in
    both.
  \end{itemize}
\item
  So the treatment may \textbf{turn off}; with absorbing treatment it
  reduces to a joiner average. Also the tool for \textbf{non-binary} and
  \textbf{continuous} treatments (\texttt{did\_multiplegt\_dyn}).
\item
  The price is \textbf{no carryover}: the outcome at \(t\) depends on
  treatment at \(t\) only --- the road is dry the moment it stops
  raining. No anticipation pointed the other way, and often harder to
  believe. Ask whether ``off'' is a return to baseline or a switch into
  something else.
\end{itemize}
\end{frame}

\begin{frame}{Imputation: Borusyak, Jaravel and Spiess (2024)}
\protect\phantomsection\label{imputation-borusyak-jaravel-and-spiess-2024}
\begin{enumerate}
\tightlist
\item
  Using \textbf{untreated cells only} --- never-treated and
  not-yet-treated --- fit
  \(Y_{it} = \alpha_i + \tau_t + \varepsilon_{it}\).
\item
  Predict \(\hat Y^0_{it} = \hat\alpha_i + \hat\tau_t\) for every
  treated cell: the unit effect from the unit's own pre-period, the year
  effect from whoever is untreated that year.
\item
  \(\hat\tau_{it} = Y_{it} - \hat Y^0_{it}\) for each treated cell;
  average with the weights your target parameter demands.
\item
  Standard errors from the implied linear combination of outcomes.
\end{enumerate}

\begin{itemize}
\tightlist
\item
  Baker data by hand, feasible cells (years to 2003; no year effect is
  estimable once everyone is treated): simple ATT \(68.33\) against CS's
  \(68.34\); cohort ATTs \(95.0\), \(52.0\), \(21.0\).
\item
  The assumption is stronger: parallel trends \textbf{in every period
  and for every unit}, not just from the baseline \(g-1\). In exchange
  it uses every pre-period and is the most efficient under
  homoskedasticity. Wooldridge (2021) is numerically the same in
  balanced panels; Gardner (2021) is the two-stage version.
\item
  Choose by the assumption, not the standard error. With serially
  correlated shocks, CS asks for less.
\end{itemize}
\end{frame}

\begin{frame}{Does it matter? Finance and accounting}
\protect\phantomsection\label{does-it-matter-finance-and-accounting}
\begin{itemize}
\tightlist
\item
  Baker, Larcker and Wang (2022) found over \(700\) DiD papers in the
  top finance and accounting journals, about \(400\) with staggered
  adoption, nearly all estimated by TWFE.
\item
  They replicated two: \textbf{bank deregulation} (interstate branching)
  and \textbf{board reforms} across countries. With the estimators of
  this section, the headline results were weakened or overturned.
\item
  The mechanism is the one on the Baker slide: effects that build over
  time, cohorts that differ, and a regression that subtracts early
  adopters' growing effects from late adopters' estimates.
\end{itemize}

\vspace{0.3em}

\begin{itemize}
\tightlist
\item
  What this does \textbf{not} say: that staggered DiD is dead, or that
  every old result is wrong. It says that for a decade the field
  reported a number without knowing which comparisons were inside it.
\item
  If you are in a staggered design: say which cells are estimable, which
  control group you use, and how you aggregate. If you report TWFE,
  report the Bacon weights beside it.
\end{itemize}
\end{frame}

\section{6. Reading an Event Study}\label{reading-an-event-study}

\begin{frame}{Long differences and short gaps (Roth 2024)}
\protect\phantomsection\label{long-differences-and-short-gaps-roth-2024}
The OLS event study compares every period to \textbf{one} baseline,
\(\ell = -1\): \[
\hat\delta_{-3} = (\bar Y_T^{\,-3} - \bar Y_T^{\,-1}) - (\bar Y_C^{\,-3} - \bar Y_C^{\,-1}) \qquad\text{(long difference)}.
\] Some packages instead compare each period to the \textbf{one before
it}: \[
\hat\delta_{-3} = (\bar Y_T^{\,-3} - \bar Y_T^{\,-2}) - (\bar Y_C^{\,-3} - \bar Y_C^{\,-2}) \qquad\text{(short gap)}.
\]

\begin{center}\small
\begin{tabular}{lrrrrrrrrrrr}
\toprule
$\ell$ & $-5$ & $-4$ & $-3$ & $-2$ & $-1$ & $0$ & $1$ & $2$ & $3$ & $4$ & $5$ \\
\midrule
Long difference & -1.2 & -0.9 & -0.6 & -0.3 & 0.0 & 3.3 & 3.6 & 3.9 & 4.2 & 4.5 & 4.8 \\
Short gap & -- & 0.3 & 0.3 & 0.3 & 0.3 & 3.3 & 0.3 & 0.3 & 0.3 & 0.3 & 0.3 \\
\bottomrule
\end{tabular}
\end{center}

\begin{itemize}
\tightlist
\item
  Noise-free example: treated group trending \(0.3\) a year faster than
  controls throughout, plus an effect of \(3\) at \(\ell = 0\). The long
  differences slope toward zero and keep climbing after treatment; the
  short gaps are a flat \(0.3\) with one spike. \textbf{Both are correct
  arithmetic. They answer different questions.}
\item
  The tell: a plotted coefficient \emph{with a confidence interval} at
  \(\ell = -1\) means short gaps.
\end{itemize}
\end{frame}

\begin{frame}[fragile]{Which to show, and what the imputation kink
means}
\protect\phantomsection\label{which-to-show-and-what-the-imputation-kink-means}
\begin{itemize}
\tightlist
\item
  Post-treatment effects are always long differences from the baseline.
  A falsification should use the \emph{same} calculation on periods
  where the effect is known to be zero --- so show pre-period
  \textbf{long differences} too. Readers trained on TWFE will read them
  that way regardless.
\item
  Stata: \texttt{csdid\ ...\ long2} (the default in \texttt{csdid2}); R:
  \texttt{did::aggte(...,\ base\_period\ =\ "universal")}. The dCDH
  Stata default is already long; check the R package.
\end{itemize}

\vspace{0.3em}

\begin{itemize}
\tightlist
\item
  \textbf{Imputation estimators look different before treatment.} BJS
  uses \emph{all} pre-periods to build the counterfactual, so the post
  coefficients assume parallel trends everywhere while the pre
  coefficients test it. The series can kink at \(\ell = -1\) even when
  the trends are right. Not a bug; a different object.
\item
  Hence the \textbf{``all the DiDs'' plot}, six estimators on one graph,
  is less reassuring than it looks: different control groups, baselines,
  weights and assumptions. They are not robustness checks for one
  another. Pick the estimator from the design, before the results, and
  say why. If you truly cannot choose ex ante, report both.
\end{itemize}
\end{frame}

\begin{frame}[fragile]{Honest DiD (Rambachan and Roth 2023)}
\protect\phantomsection\label{honest-did-rambachan-and-roth-2023}
\begin{itemize}
\tightlist
\item
  We have always used pre-trends as \emph{evidence about} parallel
  trends, informally. Make it formal: let the post-period violation be
  \textbf{at most \(M\) times the largest pre-period violation}, and
  report the set of ATTs consistent with the data for each \(M\).
\item
  \(M = 0\) is parallel trends. \(M = 1\) allows the post-period to go
  as wrong as the pre-period ever did. The confidence set widens with
  \(M\); the question is \textbf{at what \(M\) it first covers zero.}
\item
  Cunningham and Goodman-Bacon (2024), legal services and welfare
  caseloads: a pre-trend of about \(0.2\) cases per \(1{,}000\) women
  per year. Allowing post-period violations from half to twice that, the
  long-run estimates stay between \(7.5\) and \(10.5\), and the null of
  no effect is rejected at every event time but the last two.
\end{itemize}

\vspace{0.3em}

\begin{itemize}
\tightlist
\item
  Not a test that parallel trends failed. A statement of \textbf{how far
  it can fail before the conclusion does} --- a robustness check that
  removes an assumption instead of adding a model. R and Stata:
  \texttt{honestdid}.
\end{itemize}
\end{frame}

\section{7. A Checklist}\label{a-checklist}

\begin{frame}{Ten steps, in order}
\protect\phantomsection\label{ten-steps-in-order}
\begin{enumerate}
\tightlist
\item
  \textbf{Define the target parameter}: which potential outcomes, which
  population, which weights.
\item
  \textbf{Count the units in each cohort}, including never- and
  always-treated.
\item
  \textbf{Plot the rollout.}
\item
  \textbf{Pick the control group} and defend it.
\item
  \textbf{Unconditional or conditional} parallel trends; if conditional,
  which covariates.
\item
  \textbf{Check imbalance}, several ways.
\item
  \textbf{Plot average outcomes by cohort}: the first time you look at
  \(Y\).
\item
  \textbf{Choose the estimator} from the design; state its assumptions.
\item
  \textbf{Falsify}: event study, placebos, sensitivity bounds. Know how
  your pre-trends were computed.
\item
  \textbf{DDDiD}, Don't Do Diff-in-Diff. Fold a losing hand early.
\end{enumerate}

\begin{itemize}
\tightlist
\item
  Steps \(1\)--\(6\) never look at the outcome (Rubin 2008). Running
  example: Dias and Fontes (2024), Brazil's psychiatric reform.
  Community mental-health centres (CAPS) opened in \(1{,}600\)
  municipalities over \(2002\)--\(2016\); outcome, the homicide rate.
\end{itemize}
\end{frame}

\begin{frame}{Step 1: the parameter comes before the model}
\protect\phantomsection\label{step-1-the-parameter-comes-before-the-model}
Ten people in two municipalities, each with both potential outcomes for
``murdered'':

\begin{center}\scriptsize
\begin{tabular}{lcc}
\toprule
Municipality & Residents & Average $\delta_i$ \\
\midrule
1 (urban) & 8 & $0.50$ \\
2 (rural) & 2 & $-1.00$ \\
\midrule
ATE, average municipality & & $-0.25$ \\
ATE, average person & & $0.20$ \\
\bottomrule
\end{tabular}
\end{center}

\begin{itemize}
\tightlist
\item
  Same ten effects, two correct averages, opposite signs. Weighting by
  population recovers the per-person number from municipal data; not
  weighting answers for the average municipality.
\item
  Solon, Haider and Wooldridge (2015): the two diverge when effects are
  \textbf{heterogeneous} and people \textbf{sort} on them. Cunningham's
  Texas simulation: sort the largest effects into the five big counties
  and per-person stays \(2\) while per-county becomes \(-0.88\).
\item
  Neither is biased. Choose by \textbf{who decides}: a mayor cares about
  the average municipality, a ministry about the average person. Decide,
  write it down, then estimate.
\end{itemize}
\end{frame}

\begin{frame}[fragile]{Steps 2--4: count, plot, choose controls}
\protect\phantomsection\label{steps-24-count-plot-choose-controls}
\begin{center}\small
\begin{tabular}{lcc}
\toprule
CAPS cohort & Municipalities & Share of treated \\
\midrule
2003 & 47 & 0.03 \\
2006 (largest) & 216 & 0.16 \\
2016 & 78 & 0.06 \\
Treated 2003--2016 & 1,344 & 1.00 \\
Always-treated (2002), excluded & 296 & -- \\
Never-treated & 3,836 & -- \\
\bottomrule
\end{tabular}
\end{center}

\begin{itemize}
\tightlist
\item
  \textbf{Always-treated} units have no baseline: drop them. The largest
  cohort will dominate any aggregation; know that before you aggregate.
\item
  \textbf{Plot the rollout} (\texttt{panelView} in R and Stata): every
  unit, every year, treated or not. The reader sees the not-yet-treated
  pool shrink.
\item
  \textbf{Control group.} The more random the assignment, the less the
  choice matters. CAPS placement was not random; the never-treated are
  small and rural, so not-yet-treated units --- which did get a CAPS
  eventually --- may be the closer comparison. Decide \emph{before}
  results. Unsure? Run both and report both.
\end{itemize}
\end{frame}

\begin{frame}[fragile]{Steps 5--6: covariates and balance}
\protect\phantomsection\label{steps-56-covariates-and-balance}
\begin{itemize}
\tightlist
\item
  Conditional parallel trends needs the covariates that drive
  \textbf{both} CAPS placement and homicide trends: economics,
  demographics, urbanization, policing, social programmes. Cunningham
  takes Dias and Fontes's list rather than building his own: a
  specification chosen by someone who had not seen these results.
\item
  Standardized differences at baseline, treated vs.~never-treated: rural
  share \(-0.60\), population \(0.69\), health spending \(-0.45\), young
  women aged 20--29 \(0.75\). Many past \(0.25\). Those go in.
\item
  Then estimate the \textbf{propensity score} and look at it ---
  histogram, and a count of controls near one. He finds \(11\) control
  municipalities above \(0.995\); one is at \(0.999971\), an
  inverse-probability weight of \[
  \frac{0.999971}{1 - 0.999971} \approx 34{,}000 .
  \] One municipality worth thirty-four thousand. Trim before
  estimating. R's \texttt{did} trims automatically; not every package
  does.
\end{itemize}
\end{frame}

\begin{frame}[fragile]{Steps 7--8: look at the outcome, then estimate}
\protect\phantomsection\label{steps-78-look-at-the-outcome-then-estimate}
\begin{itemize}
\tightlist
\item
  Mean homicide rate by cohort over time: trends similar,
  \textbf{levels} very different, and the never-treated lowest of all.
  That is the rural pattern from the balance table: \(76\) percent of
  Brazil's homicides are urban. Kahn-Lang and Lang again --- explain the
  levels.
\item
  Estimator: Callaway--Sant'Anna, doubly robust, not-yet-treated
  controls, long differences, with the imbalanced covariates. Stata:
  \texttt{csdid2\ homicide\_rate\ \$controls,\ gvar(g)\ ivar(cod)\ time(ano)\ long2\ method(drimp)\ notyet}.
\item
  Results, all leads and lags: pre-period average \(-0.073\)
  (\(0.069\)); post average \(0.495\) (\(0.080\)), rising to about \(1\)
  extra homicide per \(10{,}000\) after eleven years.
\item
  But the far right of that plot is \textbf{three early cohorts}, and
  the far left is late adopters. Restrict to a window where most cohorts
  contribute, \(-5\) to \(+3\): effects positive from the first full
  year, about \(0.275\) per \(10{,}000\) at \(\ell = +3\) against a
  baseline of \(1.38\) --- roughly \(20\) percent. The lead at
  \(\ell = -4\) is negative and significant.
\end{itemize}
\end{frame}

\begin{frame}{Steps 9--10: how wrong can parallel trends be, and when to
quit}
\protect\phantomsection\label{steps-910-how-wrong-can-parallel-trends-be-and-when-to-quit}
\begin{itemize}
\tightlist
\item
  Honest DiD on the average of the four post lags, \(0.234\): the set of
  estimates excludes zero up to \(M = 1.5\) --- post-period violations
  one and a half times the worst pre-period violation. The result
  survives.
\item
  Cunningham was sceptical of this finding and now is not: a different
  estimator, a design checklist, and bounds that tolerate large
  violations all agree. Moving patients from hospitals to
  under-resourced community centres raised homicide.
\end{itemize}

\vspace{0.3em}

\begin{itemize}
\tightlist
\item
  \textbf{Step 10.} Every estimator in this lecture rests on parallel
  trends. When you cannot believe it with a straight face --- the
  mechanism in Ghanem, Sant'Anna and Wüthrich (2024) is clearly at work,
  support is unattainable, the pre-trends are a roller coaster and the
  bounds collapse at \(M = 0.2\) --- put the design down. Imbens's
  ``DDDiD''.
\item
  The next lecture is one of the alternatives: synthetic control.
\end{itemize}
\end{frame}

\begin{frame}{Conclusion}
\protect\phantomsection\label{conclusion}
\begin{enumerate}
\tightlist
\item
  \textbf{Sampling can break parallel trends} with no change in
  behaviour: attrition that depends on \(Y^0\), or a repeated
  cross-section whose composition shifts at treatment.
\item
  \textbf{Triple differences} assumes equal bias in two DiDs, not zero
  bias in one.
\item
  \textbf{Covariates} make parallel trends conditional. IPW, outcome
  regression and doubly robust estimators use it honestly; TWFE with
  additive controls also assumes homogeneous effects and constant
  covariate returns.
\item
  \textbf{Under differential timing, TWFE averages every
  \(2 \times 2\)}, including ones whose control was already treated.
  With dynamic effects it can get the sign wrong while every textbook
  assumption holds.
\item
  \textbf{CS, SA, dCDH and imputation} estimate group-time cells with
  clean controls and aggregate with weights you choose. They differ in
  control group, baseline and the parallel trends they need: choose by
  design.
\item
  \textbf{Read event studies carefully}: long differences vs.~short
  gaps, the imputation kink, honest bounds instead of ``the pre-trends
  look fine.''
\item
  \textbf{Checklist.} Parameter, cohorts, rollout, controls, covariates,
  balance; then the outcome, the estimator, the bounds. If it is not
  credible, walk away.
\end{enumerate}
\end{frame}

\begin{frame}{References}
\protect\phantomsection\label{references}
\tiny

\begin{itemize}
\tightlist
\item
  Abadie, A. (2005). Semiparametric Difference-in-Differences
  Estimators. \emph{Review of Economic Studies} 72(1), 1--19.
\item
  Baker, A. C., D. F. Larcker and C. C. Y. Wang (2022). How Much Should
  We Trust Staggered Difference-in-Differences Estimates? \emph{Journal
  of Financial Economics} 144(2), 370--395.
\item
  Baker, A., B. Callaway, S. Cunningham, A. Goodman-Bacon and P.
  Sant'Anna (2025). Difference-in-Differences Designs: A Practitioner's
  Guide. Working paper.
\item
  Borusyak, K., X. Jaravel and J. Spiess (2024). Revisiting Event-Study
  Designs: Robust and Efficient Estimation. \emph{Review of Economic
  Studies} 91(6), 3253--3285.
\item
  Braghieri, L., R. Levy and A. Makarin (2022). Social Media and Mental
  Health. \emph{AER} 112(11), 3660--3693.
\item
  Callaway, B. and P. Sant'Anna (2021). Difference-in-Differences with
  Multiple Time Periods. \emph{Journal of Econometrics} 225(2),
  200--230.
\item
  de Chaisemartin, C. and X. D'Haultfoeuille (2020). Two-Way Fixed
  Effects Estimators with Heterogeneous Treatment Effects. \emph{AER}
  110(9), 2964--2996.
\item
  Cunningham, S. (2025). \emph{Causal Inference: The Remix}, ch.~10 ---
  the source for this lecture, including the Baker simulation data.
  Castle-doctrine data: Cheng and Hoekstra (2013), as in Lecture 7a.
\item
  Cunningham, J. and A. Goodman-Bacon (2024). Legal Services, Divorce
  and Welfare. Working paper.
\item
  Dias, M. and L. F. Fontes (2024). The Effects of a Large-Scale Mental
  Health Reform: Evidence from Brazil. Working paper.
\item
  Goodman-Bacon, A. (2021). Difference-in-Differences with Variation in
  Treatment Timing. \emph{Journal of Econometrics} 225(2), 254--277.
\item
  Gruber, J. (1994). The Incidence of Mandated Maternity Benefits.
  \emph{AER} 84(3), 622--641.
\item
  Heckman, J., H. Ichimura and P. Todd (1997). Matching as an
  Econometric Evaluation Estimator. \emph{Review of Economic Studies}
  64(4), 605--654.
\item
  Hong, S.-H. (2013). Measuring the Effect of Napster on Recorded Music
  Sales: Difference-in-Differences Estimates under Compositional
  Changes. \emph{Journal of Applied Econometrics} 28(2), 297--324.
\item
  Kahn-Lang, A. and K. Lang (2019). The Promise and Pitfalls of
  Differences-in-Differences. \emph{JBES} 38(3), 613--620.
\item
  Olden, A. and J. Møen (2022). The Triple Difference Estimator.
  \emph{Econometrics Journal} 25(3), 531--553.
\item
  Rambachan, A. and J. Roth (2023). A More Credible Approach to Parallel
  Trends. \emph{Review of Economic Studies} 90(5), 2555--2591.
\item
  Roth, J. (2024). Interpreting Event-Studies from Recent
  Difference-in-Differences Methods. Working paper.
\item
  Sant'Anna, P. and J. Zhao (2020). Doubly Robust
  Difference-in-Differences Estimators. \emph{Journal of Econometrics}
  219(1), 101--122.
\item
  Solon, G., S. Haider and J. Wooldridge (2015). What Are We Weighting
  For? \emph{Journal of Human Resources} 50(2), 301--316.
\item
  Sun, L. and S. Abraham (2021). Estimating Dynamic Treatment Effects in
  Event Studies with Heterogeneous Treatment Effects. \emph{Journal of
  Econometrics} 225(2), 175--199.
\item
  Wooldridge, J. M. (2021). Two-Way Fixed Effects, the Two-Way Mundlak
  Regression, and Difference-in-Differences Estimators. Working paper.
\end{itemize}
\end{frame}

\section{Backup}\label{backup}

\begin{frame}[fragile]{Stata: the estimators of this lecture}
\protect\phantomsection\label{stata-the-estimators-of-this-lecture}
\begin{Shaded}
\begin{Highlighting}[]
\KeywordTok{use}\NormalTok{ https:}\CommentTok{//github.com/scunning1975/mixtape/raw/master/baker.dta, clear}
\KeywordTok{areg} \FunctionTok{y}\NormalTok{ i.}\FunctionTok{year}\NormalTok{ treat, a(id) }\KeywordTok{robust}                   \CommentTok{// TWFE static = {-}6.69}
\NormalTok{bacondecomp }\FunctionTok{y}\NormalTok{ treat, ddetail                         }\CommentTok{// Goodman{-}Bacon weights}
\NormalTok{csdid2 }\FunctionTok{y}\NormalTok{, ivar(id) time(}\FunctionTok{year}\NormalTok{) gvar(treat\_date) notyet method(stdipw) long2}
\KeywordTok{estat}\NormalTok{ simple                                         }\CommentTok{// 68.3}
\KeywordTok{estat} \FunctionTok{group}                                          \CommentTok{// 56}
\KeywordTok{estat}\NormalTok{ event                                          }\CommentTok{// relative{-}time aggregation}
\NormalTok{eventstudyinteract }\FunctionTok{y}\NormalTok{ lead* lag*, cohort(treat\_date) control\_cohort(}\FunctionTok{last}\NormalTok{) absorb(id }\FunctionTok{year}\NormalTok{)}
\NormalTok{did\_imputation }\FunctionTok{y}\NormalTok{ id }\FunctionTok{year}\NormalTok{ treat\_date, autosample      }\CommentTok{// Borusyak{-}Jaravel{-}Spiess}
\NormalTok{did\_multiplegt\_dyn }\FunctionTok{y}\NormalTok{ id }\FunctionTok{year}\NormalTok{ treat                   }\CommentTok{// de Chaisemartin{-}D\textquotesingle{}Haultfoeuille}
\NormalTok{honestdid, pre(1/4) }\KeywordTok{post}\NormalTok{(5/8) mvec(0(0.5)2)          }\CommentTok{// Rambachan{-}Roth, after an event study}
\NormalTok{drdid }\FunctionTok{y}\NormalTok{ x1 }\KeywordTok{x2}\NormalTok{, time(}\FunctionTok{year}\NormalTok{) ivar(id) tr(treat) }\OtherTok{all}     \CommentTok{// IPW, OR, DR with covariates}
\end{Highlighting}
\end{Shaded}

\begin{itemize}
\tightlist
\item
  R equivalents: \texttt{bacondecomp}, \texttt{did} (\texttt{att\_gt},
  \texttt{aggte}), \texttt{fixest::sunab}, \texttt{didimputation},
  \texttt{DIDmultiplegtDYN}, \texttt{HonestDiD}, \texttt{DRDID}. This
  deck computed CS, Bacon and the imputation estimator from group means
  with \texttt{fixest} and base R only.
\end{itemize}
\end{frame}




\end{document}
